\section{Question 1}

\subsection{Question 1.1}

Input Feature Vector: 
\[ x = [\text{account\_age, monthly\_bill, support\_calls, contract}]\]
So: 
\[x \in \mathbb{R}^4\]

Target Variable:
\[y = \text{churn}\]

where 
\[y = 
\begin{cases} 
1, & \text{if customer churned} \\
0, & \text{if customer did not churn} 
\end{cases}\]


This is a classification problem because the task at hand is predicting whether a customer stayed or 
left the telecommunications company. Therefore, it is a binary classification problem since there are
only two possible outcomes (not predicting a continuous value which regression is used for).

\subsection{Question 1.2}

The churn rate of the dataset is 0.12 or 12\%. Therefore, the dataset is an imbalanced dataset since
only 12\% of the 200 customers churned (left the company).

The pandas expression to compute the mean monthly bill separately for churners and non-churners is given as follows:
\[\text{data.groupby('churn')['monthly\_bill'].mean()}\]
which outputs the following:

\begin{table}[h]
    \centering
    \caption{Mean Value by Churn Status}
    \label{tab:mean_by_churn}
    \begin{tabular}{cc}
        \hline
        \textbf{Churn} & \textbf{Mean} \\
        \hline
        \text{non-churners} & 109.10 \\
        \text{churners} & 104.07 \\
        \hline
    \end{tabular}
\end{table}

\subsection{Question 1.3}

Because raw accuracy can give a false sense of the model's performance. For example if the
churn percentage is only 15\%, a model that predicts everyone does not churn achieves a
raw accuracy of 85\%. This is clearly misleading because the classifier failed to predict 
any positive cases (churns) which are the customers the company wants to identify so that
actions can be taken to keep them or future similar customers to stay.

A metric that handles class imbalance well is Recall. Recall is a measure
of the true positive predictions across all the actual positive cases. In this scenario,
it measures how many actual churners the model correctly identified. This is useful because 
missing a customer that is likely to leave is more costly can incorrectly flagging a customer
who would have actually stayed.
